% Supporting Information (separate file, hosted online without typesetting). Same macros and generated tables.
\documentclass[12pt]{article}
\usepackage[margin=1in]{geometry}\usepackage{mathptmx}\usepackage{setspace}\onehalfspacing
\usepackage{booktabs}\usepackage{array}\usepackage{longtable}\usepackage{amsmath}\usepackage{graphicx}\usepackage{rotating}\usepackage{xr-hyper}\externaldocument{main}\usepackage[hidelinks]{hyperref}\usepackage{apacite}
% generated by code/make_tables_cgir.py on 2026-10-28 — every macro traces to a results file (see the generator for keys)
\newcommand{\NOldN}{464}
\newcommand{\NOldDisposed}{257}
\newcommand{\NOldRetained}{89}
\newcommand{\NOldOther}{43}
\newcommand{\NOldUnresolved}{75}
\newcommand{\NOldPriced}{380}
\newcommand{\NOldPricedOneY}{366}
\newcommand{\NOldOwnership}{387}
\newcommand{\NNewN}{294}
\newcommand{\NNewDisposed}{145}
\newcommand{\NNewRetained}{90}
\newcommand{\NNewOther}{29}
\newcommand{\NNewUnresolved}{30}
\newcommand{\NNewPriced}{260}
\newcommand{\NNewPricedOneY}{249}
\newcommand{\NNewOwnership}{283}
\newcommand{\NAllN}{758}
\newcommand{\NAllDisposed}{402}
\newcommand{\NAllRetained}{179}
\newcommand{\NAllOther}{72}
\newcommand{\NAllUnresolved}{105}
\newcommand{\NAllPriced}{640}
\newcommand{\NAllPricedOneY}{615}
\newcommand{\NAllOwnership}{670}
\newcommand{\NScoreOneY}{559}
\newcommand{\NLabelOneY}{471}
\newcommand{\NGateOne}{579}
\newcommand{\AUCPooledLLM}{0.761}
\newcommand{\AUCPooledFin}{0.587}
\newcommand{\AUCNewLLM}{0.813}
\newcommand{\AUCNewFin}{0.577}
\newcommand{\DeltaPooled}{0.174}
\newcommand{\DeltaPooledLo}{0.117}
\newcommand{\DeltaPooledHi}{0.235}
\newcommand{\AUCOldIdent}{0.706}
\newcommand{\AUCOldAnon}{0.691}
\newcommand{\AUCNewIdent}{0.810}
\newcommand{\AUCNewAnon}{0.803}
\newcommand{\ClaudeOutsideKnowledgeOld}{43.0}
\newcommand{\AUCClaudeFlaggedAnonOld}{0.763}
\newcommand{\AUCClaudeFlaggedIdentOld}{0.764}
\newcommand{\DisposedSharePooled}{69.3}
\newcommand{\GapOldLo}{$-$0.006}
\newcommand{\GapOldHi}{0.039}
\newcommand{\AUCEarlyPooled}{0.780}
\newcommand{\AUCEarlyNew}{0.919}
\newcommand{\AUCEarlyOld}{0.698}
\newcommand{\AUCLateVsNever}{0.708}
\newcommand{\SpearmanAnon}{$-$0.135}
\newcommand{\NAUCEarlyAnon}{339}
\newcommand{\NAUCEarly}{412}
\newcommand{\SpearmanAnonP}{0.001}
\newcommand{\RankAnonCoef}{$-$0.195}
\newcommand{\RankAnonP}{0.005}
\newcommand{\AUCEarlyAnon}{0.817}
\newcommand{\QQDiff}{$-$4.4}
\newcommand{\QQLo}{$-$22.7}
\newcommand{\QQHi}{+17.3}
\newcommand{\QQWithinYear}{$-$11.7}
\newcommand{\QQWithinYearP}{0.154}
\newcommand{\SpearmanScoreBhar}{$-$0.136}
\newcommand{\SpearmanScoreBharP}{0.001}
\newcommand{\SpearmanWithinYear}{$-$0.134}
\newcommand{\SpearmanWithinYearP}{0.002}
\newcommand{\AnnCarSlope}{0.037}
\newcommand{\AnnCarSlopeSE}{0.019}
\newcommand{\AnnCarSlopeP}{0.046}
\newcommand{\CTHiLo}{$-$12.2}
\newcommand{\CTHiLoT}{$-$1.68}
\newcommand{\FFHiLo}{$-$0.52}
\newcommand{\FFHiLoT}{$-$0.96}
\newcommand{\DriftMean}{$-$6.0}
\newcommand{\DriftMedian}{$-$15.9}
\newcommand{\DriftNeg}{66}
\newcommand{\DriftLo}{$-$11.6}
\newcommand{\DriftHi}{$-$0.1}
\newcommand{\AnnCarMean}{+0.58}
\newcommand{\AnnCarT}{1.30}
\newcommand{\AnnCarN}{583}
\newcommand{\FFAllAlpha}{$-$0.02}
\newcommand{\FFAllT}{$-$0.07}
\newcommand{\EarlyN}{195}
\newcommand{\EarlyMean}{$-$14.2}
\newcommand{\EarlyMedian}{$-$26.6}
\newcommand{\NeverN}{141}
\newcommand{\NeverMean}{+1.4}
\newcommand{\NeverMedian}{$-$8.7}
\newcommand{\EarlyNeverDiff}{$-$15.6}
\newcommand{\EarlyNeverLo}{$-$34.2}
\newcommand{\EarlyNeverHi}{+1.3}
\newcommand{\EarlyNeverWelchP}{0.082}
\newcommand{\EarlyNeverMWP}{0.001}
\newcommand{\LateMean}{$-$3.8}
\newcommand{\LateNeverP}{0.626}
\newcommand{\DispRetDiff}{$-$10.3}
\newcommand{\DispRetMWP}{0.010}
\newcommand{\DispRetWelchP}{0.219}
\newcommand{\RankSpecOneCoef}{$-$0.191}
\newcommand{\RankSpecOneP}{0.004}
\newcommand{\RankSpecOneClP}{0.001}
\newcommand{\RankSpecTwoCoef}{$-$0.192}
\newcommand{\RankSpecTwoP}{0.003}
\newcommand{\RankSpecTwoClP}{$<$0.001}
\newcommand{\RankSpecThreeCoef}{$-$0.195}
\newcommand{\RankSpecThreeP}{0.004}
\newcommand{\RankSpecThreeClP}{$<$0.001}
\newcommand{\RankSpecThreeBCoef}{$-$0.229}
\newcommand{\RankSpecThreeBP}{0.019}
\newcommand{\RankSpecThreeBClP}{0.002}
\newcommand{\RankSpecThreeCCoef}{$-$0.168}
\newcommand{\RankSpecThreeCP}{0.097}
\newcommand{\RankSpecThreeCClP}{0.047}
\newcommand{\RankSpecFourCoef}{$-$0.207}
\newcommand{\RankSpecFourP}{0.047}
\newcommand{\RankSpecFourClP}{0.042}
\newcommand{\RankScoreCoefMin}{$-$0.20}
\newcommand{\RankScoreCoefMax}{$-$0.19}
\newcommand{\JointEarlyCoef}{$-$0.151}
\newcommand{\JointEarlyP}{0.083}
\newcommand{\JointScoreCoef}{0.135}
\newcommand{\JointScoreP}{0.549}
\newcommand{\BlockMedian}{34.6}
\newcommand{\LargestMedian}{21.8}
\newcommand{\NOwnership}{551}
\newcommand{\ScoreOnBlock}{$-$0.113}
\newcommand{\ScoreOnBlockP}{0.013}
\newcommand{\ScoreOnBlockClP}{0.004}
\newcommand{\LowBlockDisp}{74}
\newcommand{\HighBlockDisp}{62}
\newcommand{\LowBlockEarly}{52}
\newcommand{\HighBlockEarly}{43}
\newcommand{\LowBlockMedian}{$-$26.9}
\newcommand{\HighBlockMedian}{$-$9.1}
\newcommand{\LowBlockMean}{$-$13.7}
\newcommand{\HighBlockMean}{$-$0.5}
\newcommand{\LowHighMWP}{$<$0.001}
\newcommand{\LowHighDiff}{$-$13.1}
\newcommand{\LowBlockMax}{27}
\newcommand{\HighBlockMin}{43}
\newcommand{\RankOnBlock}{0.14}
\newcommand{\RankOnBlockP}{0.062}
\newcommand{\ProtectPre}{0.5}
\newcommand{\ProtectPost}{95.5}
\newcommand{\ListingPre}{9.1}
\newcommand{\ListingPost}{77.3}
\newcommand{\CarPre}{+0.06}
\newcommand{\CarPost}{+5.29}
\newcommand{\CarPrePostP}{$<$0.001}
\newcommand{\ScorePre}{0.528}
\newcommand{\ScorePost}{0.510}
\newcommand{\ScorePrePostP}{0.616}
\newcommand{\EarlyPre}{40.2}
\newcommand{\NPostOutcomes}{19}
\newcommand{\EarlyPost}{36.8}
\newcommand{\EarlyPrePostP}{1.000}
\newcommand{\NPost}{49}
\newcommand{\NPostFilings}{44}
\newcommand{\FilingsPerYearPre}{39.0}
\newcommand{\FilingsPerYearPost}{16.4}
\newcommand{\NDefs}{10}
\newcommand{\IPOShare}{3.3}
\newcommand{\NClean}{581}
\newcommand{\KLAnnCar}{0.853}
\newcommand{\KLBhar}{$-$10.995}
\newcommand{\KLCar}{$-$7.857}
\newcommand{\KLDisposed}{69.3}
\newcommand{\KLN}{525}
\newcommand{\KLOwnMedian}{14.71}
\newcommand{\KLOwnMean}{18.432}
\newcommand{\RepeatCorrA}{0.96}
\newcommand{\RepeatCorrB}{0.99}
\newcommand{\BetweenModelCorr}{0.87}
\newcommand{\NRepeat}{245}
\newcommand{\KRXCompustatCorr}{0.9998}
\newcommand{\NPriceOverlap}{233}
\newcommand{\GlitchExcluded}{14}
\newcommand{\OwnPairs}{357}
\newcommand{\OwnWithinOne}{354}
\newcommand{\OwnDisagree}{3}
\newcommand{\AUCTemporalFin}{0.611}
\newcommand{\AUCTemporalLLM}{0.813}
\newcommand{\AnnCarSlopePP}{3.7}
\newcommand{\RankSpecThreePctl}{19.5}
\newcommand{\GPTOutsideKnowledgeOld}{3.8}
\newcommand{\CleanDisposedShare}{69.2}
\newcommand{\KLListed}{4}
\newcommand{\BlockOnDisposed}{$-$0.092}
\newcommand{\BlockOnDisposedSE}{0.129}
\newcommand{\BlockOnEarlyAlone}{$-$0.141}
\newcommand{\BlockOnEarlyAloneSE}{0.144}
\newcommand{\BlockOnEarlyWithScore}{0.113}
\newcommand{\BlockOnEarlyWithScoreSE}{0.133}
\newcommand{\ScoreOnEarlyWithBlock}{1.235}
\newcommand{\BlockOnCar}{0.030}
\newcommand{\BlockOnCarSE}{0.023}
\newcommand{\EarlyCoefBhar}{$-$0.243}
\newcommand{\LowBlockCoefBhar}{$-$0.231}
\newcommand{\InterCoefBhar}{0.223}
\newcommand{\InterPBhar}{0.114}
\newcommand{\EarlyCoefRank}{$-$0.136}
\newcommand{\LowBlockCoefRank}{$-$0.111}
\newcommand{\InterCoefRank}{0.114}
\newcommand{\InterPRank}{0.071}
\newcommand{\EarlyInLowBlockRank}{$-$0.022}
\newcommand{\LGScore}{0.76}
\newcommand{\LGAnnCar}{$-$7.3}
\newcommand{\LGBhar}{$-$15.3}
\newcommand{\KLKrxTotal}{948}
\newcommand{\KLKrxCaptive}{719}
\newcommand{\NScored}{720}
\newcommand{\QOneMedian}{$-$3.0}
\newcommand{\QFiveMedian}{$-$23.0}
\newcommand{\QQMedianDiff}{$-$20.0}
\newcommand{\QQMWP}{0.008}
\newcommand{\KLLGDrop}{30}
\newcommand{\NSensitive}{11}
\newcommand{\NMixed}{4}
\newcommand{\ExDoosanScore}{0.98}
\newcommand{\ExDoosanOutcome}{disposed}
\newcommand{\ExDBScore}{0.24}
\newcommand{\ExDBOutcome}{retained}
\newcommand{\ExTvingScore}{0.85}
\newcommand{\ExTvingOutcome}{disposed}
\newcommand{\NUnread}{38}
\newcommand{\NUnreadTooRecent}{15}
\newcommand{\NUnreadUnmatched}{13}
\newcommand{\NUnreadMergedLiq}{9}
\newcommand{\NUnreadClean}{1}
\newcommand{\NPostAll}{66}
\newcommand{\KLBharOne}{$-$11.0}
\newcommand{\KLOwnMeanOne}{18.4}
\newcommand{\KLOwnMedianOne}{14.7}
\newcommand{\MDEEarlyPP}{34}
\newcommand{\MDEScore}{0.10}
\newcommand{\ScoreSDPre}{0.22}

\newcommand{\tabnote}[1]{\par\vspace{4pt}\begin{minipage}{\linewidth}\singlespacing\small\emph{Note.} #1\end{minipage}}
\title{Supporting Information}\author{}\date{}
\begin{document}
\maketitle
\section{Definitions in full}
{\footnotesize% generated by code/make_tables_cgir.py on 2026-10-28 — source: DEFS list in code/make_tables_cgir.py (same rows as T_defs_short)
\begin{longtable}{>{\raggedright\arraybackslash}p{0.7cm}>{\raggedright\arraybackslash}p{2.4cm}>{\raggedright\arraybackslash}p{3.5cm}>{\raggedright\arraybackslash}p{4.8cm}>{\raggedright\arraybackslash}p{2.1cm}}
\caption{Definitions of the measures used in this paper (full form)}\label{tab:s_defs_full}\\
\toprule No. & Term & Conceptual definition & Operational rule and source & Applied in \\ \midrule \endfirsthead
\multicolumn{5}{l}{\small Table~\thetable\ (continued)}\\ \toprule No. & Term & Conceptual definition & Operational rule and source & Applied in \\ \midrule \endhead
\midrule\multicolumn{5}{r}{\small continued on next page}\\ \endfoot
\bottomrule \endlastfoot
D1 & Captive spin-off (\emph{mulchŏk bunhal}) & A corporate division under the Korean Commercial Act's division provisions in which the surviving company holds 100\% of the newly created company \protect\cite{kang2025}; unlike a US spin-off no shares are distributed to the parent's shareholders, and unlike an equity carve-out no shares are sold or listed at the division. & DART division-resolution filings (\emph{hoesa bunhal gyeoljeong}; from 2017 the material-event report) whose division type is the captive form (a mixed captive--equity division is kept when a wholly owned company is created); author rule, Section~\ref{sec:sample}. & Table~\ref{tab:sample}; all tests \\ \addlinespace[2pt]
D2 & Disposal & The event through which the value of the carved-out assets leaves the parent's minority shareholders \protect\cite{kang2025}: a sale, a listing, or a dilution of the parent's stake. & Ownership of the new company read from the parent's annual reports (the schedule of equity investments in other corporations and the notes on subsidiaries) for each fiscal year from the division to 2025, with KIND listing records; majority sale, listing, or stake below 99.5\% = disposed; 100\% throughout = retained; mergers back into the parent and liquidations excluded; author rule, Section~\ref{sec:sample}. & Tables~\ref{tab:sample}, \ref{tab:gate1}, \ref{tab:returns_by_outcome} \\ \addlinespace[2pt]
D3 & Early disposal & Disposal within the one-year return window. & first\_change\_fy $-$ division year $\le 1$ from the annual-report labels; author rule. & Tables~\ref{tab:returns_by_outcome}, \ref{tab:realized_controls}, \ref{tab:block_terciles} \\ \addlinespace[2pt]
D4 & Disposal signal & What a careful reader could infer about the controlling owner's intention from the filing alone at the announcement (\protect\citeNP{kang2025}, Section 4, call this unlikely). & Each model reads the division-resolution filing as of its date and returns a probability (0--100) that the parent will not retain 100\% within three to five fiscal years; mean over runs per model, then the two-model mean; identified condition unless stated; Section~\ref{sec:signal}. & Tables~\ref{tab:gate1}, \ref{tab:returns_by_score}, \ref{tab:controls} \\ \addlinespace[2pt]
D5 & Anonymized condition & A reading that cannot use knowledge of the named company's later history \protect\cite{glasserman2024}. & Regular-expression masking of names, group names and dates before the model reads the text; Section~\ref{sec:signal}. & Table~\ref{tab:gate1} \\ \addlinespace[2pt]
D6 & Controlling block & The controlling owner's cash-flow stake in the parent, the quantity that determines how much of a minority loss the controller bears \protect\cite{young2008}. & 2016+ announcements: the largest-shareholder table of the annual report via the DART API, total row minus treasury-share rows; earlier: the same table read from the annual report of the pre-division fiscal year (year-end column), or of the division year (year-start column) when the former is not available; Section~\ref{sec:returns}. & Tables~\ref{tab:block_terciles}, \ref{tab:block_regressions} \\ \addlinespace[2pt]
D7 & Announcement CAR & The immediate price reaction to the filing. & Market model estimated on days $-250$ to $-31$ against the KOSPI or KOSDAQ index of the parent's market; day 0 = first trading day on or after the receipt date of the original filing; Section~\ref{sec:returns}. & Tables~\ref{tab:returns_by_score}, \ref{tab:controls}, \ref{tab:reform} \\ \addlinespace[2pt]
D8 & One-year BHAR & The return to a minority holder who kept the shares after the announcement reaction. & Daily adjusted prices (KRX for 2017+; Compustat Global for 2001--2016); windows with fewer than 90\% of days are dropped; Section~\ref{sec:returns}. & Tables~\ref{tab:returns_by_score}--\ref{tab:block_regressions} \\ \addlinespace[2pt]
D9 & Within-year rank & A skew-robust ordering of outcomes within the same calendar environment. & Rank within announcement year; author rule. & Tables~\ref{tab:controls}, \ref{tab:realized_controls}, \ref{tab:block_regressions} \\ \addlinespace[2pt]
D10 & Post-reform period & Filings written under the revised disclosure form that requires the division's purpose, expected effects and shareholder-protection measures (FSC, 5 September 2022). & Receipt date $\ge$ 2022-10-01; the first filing after the cutoff that states protection measures is dated 17 October 2022; Section~\ref{sec:reform}. & Table~\ref{tab:reform} \\ \addlinespace[2pt]
\end{longtable}
}
\section{Label definitions and sensitivity}
\begin{table}[htbp]\centering\caption{Headline results with the \NSensitive{} definition-sensitive labels flipped or dropped}\label{tab:s_labels}
{\footnotesize\setlength{\tabcolsep}{3pt}% generated by code/make_tables_cgir.py on 2026-10-26 — source: label_sensitivity_20261025.json
\begin{tabular}{>{\raggedright\arraybackslash}p{6.2cm}rrr}
\toprule
 & Labels as used & Eleven set to retained & Eleven dropped \\
\midrule
Clean-label events with returns & 476 & 476 & 468 \\
Disposed share (\%) & 69.5 & 67.9 & 69.0 \\
AUC, two-model disposal score (pooled; 2017+) & 0.785; 0.830 & 0.790; 0.844 & 0.792; 0.850 \\
Disposed $-$ retained, one-year BHAR (\%) [MW $p$] & $-$10.3 [0.010] & $-$14.6 [0.002] & $-$11.8 [0.006] \\
Early $-$ never, one-year BHAR (\%) [MW $p$] & $-$15.6 [0.001] & $-$18.3 [$<$0.001] & $-$15.6 [0.001] \\
Rank of BHAR on early disposal, controls [HC1 $p$] & $-$0.079 [0.013] & $-$0.079 [0.013] & $-$0.079 [0.013] \\
Disposal rate, low / high block tercile (\%) & 74 / 62 & 72 / 59 & 74 / 61 \\
\bottomrule
\end{tabular}
}\tabnote{\NSensitive{} of the 2017--2026 labels are sensitive to the outcome definition or to a disagreement between the equity-investment schedule and the annual-report text (categories: intragroup transfer, temporary external dilution, related-party preferred shares, dilution below materiality, unresolved conflict, one tracker mismatch). Column 2 recodes all eleven as retained, the direction least favourable to the paper's results; column 3 drops them. Statistics as defined in the main text.}\end{table}
\section{Price sources}
\begin{table}[htbp]\centering\caption{Two price sources: cross-check and screen}\label{tab:s_prices}
{\footnotesize\setlength{\tabcolsep}{3pt}% generated by code/make_tables_cgir.py on 2026-10-26 — source: price_coverage_compustat_20261024.json; bhar_pooled_glitch_screen_20261025.json
\begin{tabular}{>{\raggedright\arraybackslash}p{11cm}r}
\toprule
Price-source cross-check and screen & \\
\midrule
2017+ events priced in both sources & 233 \\
Trading days compared & 110,420 \\
Median correlation of daily returns & 0.9998 \\
Share of days on which the two sources differ by more than 1 pp (\%) & 0.2 \\
Share of events with correlation below 0.95 (\%) & 2.1 \\
\midrule
2001--2016 priced events & 380 \\
\quad with a day beyond the exchange price limit in [$-$1, +250] & 26 \\
\quad of which in delisting proceedings (kept) & 12 \\
\quad suspected adjustment errors (excluded) & 14 \\
\bottomrule
\end{tabular}
}\tabnote{Daily adjusted prices from the Korea Exchange (2017--2026) and Compustat Global (2001--2016, used for 2017--2026 only in this cross-check). Returns are compared on the events priced in both sources over trading days $-250$ to $+250$. The exchange's adjusted close excludes cash dividends while Compustat's adjustment factor includes them, so small differences on ex-dividend days are expected. The screen flags a 2001--2016 event when a daily return in $[-1, +250]$ exceeds the exchange's daily price limit (15 percent before 15 June 2015, 30 percent after) by more than one percentage point; flagged events whose parent was in delisting proceedings are kept, the rest are excluded from the main results and included in a sensitivity run.}\end{table}
\section{Ownership sources}
\begin{table}[htbp]\centering\caption{Pre-announcement ownership: two sources and their agreement}\label{tab:s_ownership}
{\footnotesize\setlength{\tabcolsep}{3pt}% generated by code/make_tables_cgir.py on 2026-10-28 — source: largest_shareholder_coverage_20261025.json; ownership_llm_agreement_20261025.json
\begin{tabular}{>{\raggedright\arraybackslash}p{6.6cm}>{\raggedright\arraybackslash}p{3.2cm}>{\raggedright\arraybackslash}p{4.2cm}}
\toprule
Pre-announcement ownership of the parent & DART table (2016+) & Annual-report extraction (pre-2016) \\
\midrule
Events with a usable record & 326 & 373 passages; 357 with figures \\
Median controlling block (\%) & 36.0 & see the main text (block terciles) \\
Two-extractor agreement: largest holder within 0.5 pp & -- & 355 / 357 \\
Two-extractor agreement: block within 1 pp & -- & 354 / 357 \\
Records excluded for internal or cross-extractor disagreement & 5 & 3 \\
\bottomrule
\end{tabular}
}\tabnote{The controlling block is the largest shareholder's and related parties' common shares, excluding treasury shares, at the fiscal year-end before the announcement year. For announcements from 2016 the figure is read from the largest-shareholder table returned by the regulator's interface; for earlier announcements the same table is read from the annual report's text by two language models independently, and a record is kept when the two agree within one percentage point.}\end{table}
\section{Results by period}
\begin{table}[htbp]\centering\caption{Headline return results by period}\label{tab:s_period}
{\footnotesize\setlength{\tabcolsep}{3pt}% generated by code/make_tables_cgir.py on 2026-10-26 — source: bhar_predicted_disposal_pooled_20261025.json by_period
\begin{tabular}{lrrrr}
\toprule
Period & $n$ & Q5 $-$ Q1 BHAR (\%) [CI] & Spearman $\rho$ [$p$] & CAR slope [$p$] \\
\midrule
2001--2016 & 350 & $-$14.4 [$-$33.0, +6.6] & $-$0.177 [0.001] & 0.0588 [0.070] \\
2017--2026 & 209 & $-$5.6 [$-$29.0, +14.5] & $-$0.074 [0.290] & 0.0192 [0.361] \\
\bottomrule
\end{tabular}
}\tabnote{Listed parents with a disposal signal and a full one-year window, by announcement period. Q5 $-$ Q1 is the mean one-year market-adjusted BHAR difference between the top and bottom quintiles of the signal with a bootstrap 95 percent interval; Spearman correlation between the signal and the one-year BHAR; CAR slope is the OLS coefficient of the three-day announcement CAR on the signal with year fixed effects (HC1 $p$).}\end{table}
\bibliographystyle{apacite}
\bibliography{refs,refs_manual}
\end{document}
